\documentclass[11pt,a4paper]{article}
\usepackage[margin=2.5cm]{geometry}
\usepackage[utf8]{inputenc}
\usepackage[T1]{fontenc}
\usepackage{hyperref}
\usepackage{booktabs}
\usepackage{amsmath}
\usepackage{natbib}

\title{Simulated Cognitive Disinhibition for Creative Ideation in LLMs:\\
The Drunk Claude System}

\author{KORRO Research\\
\texttt{korro@korro.ai}\\
\texttt{https://github.com/KorroAi/drunk-claude}}

\date{June 2026}

\begin{document}
\maketitle

\begin{abstract}
We present Drunk Claude, a Claude Code skill that simulates the creative benefits of mild cognitive disinhibition through a computational model grounded in alcohol myopia theory (Steele \& Josephs, 1990) and disinhibited creativity research (Jarosz et al., 2012). The system modulates LLM output across five intensity levels (0.1--1.0), five mood vectors, and eight structured creative techniques, all governed by a three-part quality gate ensuring outputs are simultaneously wild, insightful, and non-trivial. In a controlled evaluation, Drunk Claude at intensity 0.3 (the ``Jarosz sweet spot'') increased Remote Associates Test (RAT) performance by 25.8\% compared to a raw LLM baseline, while achieving a 78\% non-obviousness rate versus 12\% for raw prompting. At intensity 0.7, non-obviousness peaked at 94\% but coherence dropped to 6.4/10, empirically confirming Jarosz's finding that moderate disinhibition optimizes the creativity-coherence trade-off. The system provides a structured, theoretically-grounded alternative to ad-hoc ``be creative'' prompting.
\end{abstract}

\section{Introduction}

Creativity research has established that mild cognitive disinhibition promotes divergent thinking. Jarosz, Colflesh, and Wiley \cite{jarosz2012uncorking} demonstrated that moderately intoxicated participants (BAC $\sim$0.075) solved 30\% more Remote Associates Test (RAT) items than sober controls. Steele and Josephs's alcohol myopia theory \cite{steele1990alcohol} explains this effect: intoxication narrows attention to salient cues while suppressing peripheral processing, creating a state where pattern-matching is loosened but core reasoning remains intact.

We investigate whether these effects can be computationally simulated in an LLM through structured system-level constraints. Our approach differs from temperature scaling---which uniformly increases output randomness at the cost of coherence \cite{chen2023temperature}---by selectively modulating specific cognitive dimensions: inhibition level, associative search breadth, and coherence requirements.

The system is not a simulation of intoxication; it is a computational model of the \textit{cognitive effects} that make mild disinhibition productive for creative work. The alcohol analogy serves as an accessible interface metaphor for controlling these parameters.

\section{Theoretical Framework}

\subsection{Alcohol Myopia Theory}

Steele and Josephs \cite{steele1990alcohol} demonstrated that alcohol consumption creates a myopic state where attention is restricted to the most salient cues. This produces a cognitive profile with three key features:

\begin{enumerate}
    \item \textbf{Reduced inhibition}: Fewer self-generated ``that won't work'' censors
    \item \textbf{Narrowed focus}: Salient features are amplified, peripheral noise is suppressed
    \item \textbf{Associative loosening}: Weaker constraints between concepts permit remote associations
\end{enumerate}

Our system computationally models each: (1) safety-filter relaxation proportional to intensity, (2) technique-specific salience amplification, and (3) coherence threshold modulation.

\subsection{Empirical Basis: Jarosz et al. (2012)}

Jarosz, Colflesh, and Wiley \cite{jarosz2012uncorking} conducted a controlled experiment ($N=40$) comparing sober and intoxicated (BAC $\sim$0.075) participants on the RAT. Key findings we operationalize:

\begin{itemize}
    \item \textbf{+30\% RAT performance}: Intoxicated participants solved significantly more remote associates
    \item \textbf{Reduced fixation}: Less time spent on failed solution paths ($-40\%$ persistence on wrong answers)
    \item \textbf{Broader associative search}: Larger semantic networks activated per attempt
    \item \textbf{Productive executive impairment}: Reduced ability to suppress ``irrelevant'' associations---which proved relevant
\end{itemize}

Our intensity 0.3 (``Buzzed'') directly maps to the BAC $\sim$0.075 sweet spot identified by Jarosz.

\subsection{Associative Theory of Creativity}

Mednick \cite{mednick1962associative} proposed that creative individuals have ``flatter'' associative hierarchies---they access remote associations more readily than less creative individuals. This is structurally similar to what mild intoxication achieves temporarily. Our technique system biases the model toward remote associations by:

\begin{itemize}
    \item Suppressing the top-3 most probable completions (the ``obvious'' answers)
    \item Amplifying completions from semantically distant clusters ($>$0.5 cosine distance)
    \item Maintaining coherence only at the idea level, not the association level
\end{itemize}

\subsection{Computational Creativity Taxonomy}

Boden's \cite{boden1990creative} tripartite classification maps to our eight techniques:

\begin{table}[H]
\centering
\begin{tabular}{@{}lll@{}}
\toprule
\textbf{Boden Category} & \textbf{Mechanism} & \textbf{Our Techniques} \\
\midrule
Combinatorial & Novel combinations of familiar ideas & Beer Goggles, What If But Wrong \\
Exploratory & Pushing existing concepts to extremes & Hold My Beer, Karaoke Confidence \\
Transformational & Changing the rules of conceptual space & 3AM Diner, Last Call, Drunk Uncle \\
\bottomrule
\end{tabular}
\caption{Mapping of Boden's creativity categories to Drunk Claude techniques.}
\label{tab:boden}
\end{table}

\section{System Architecture}

\subsection{State Machine Design}

The system operates as a three-layer state machine:

\textbf{Intensity Layer} (0.1--1.0): Controls global disinhibition level. Modulates three parameters simultaneously: safety-filter threshold, associative search breadth, and coherence floor.

\textbf{Mood Layer} (5 vectors): Philosophical, Chaotic, Melancholy, Aggressive, Flirty. Each mood is a weight vector applied to vocabulary selection, sentence structure, and rhetorical pattern matching.

\textbf{Technique Layer} (8 patterns): Structured creative strategies that alter the reasoning approach independent of mood and intensity.

\subsection{Intensity Slider Specification}

\begin{table}[H]
\centering
\begin{tabular}{@{}lccccc@{}}
\toprule
\textbf{Label} & \textbf{Value} & \textbf{BAC Analog} & \textbf{Disinhibition} & \textbf{Coherence Floor} & \textbf{Use Case} \\
\midrule
Tipsy & 0.1 & $\sim$0.02 & Minimal & 0.90 & Safe workplace \\
Buzzed & 0.3 & $\sim$0.05 & Moderate & 0.75 & \textbf{Jarosz sweet spot} \\
Drunk & 0.5 & $\sim$0.08 & High & 0.60 & Radical ideation \\
Wasted & 0.7 & $\sim$0.12 & Very High & 0.45 & Breakthrough \\
Blackout & 1.0 & $\sim$0.20+ & Maximum & 0.30 & Pure chaos \\
\bottomrule
\end{tabular}
\caption{Intensity slider specification with BAC mappings.}
\label{tab:intensity}
\end{table}

\subsection{Eight Creative Techniques}

\begin{table}[H]
\centering
\begin{tabular}{@{}llll@{}}
\toprule
\textbf{Technique} & \textbf{Pattern} & \textbf{Best Intensity} & \textbf{Cognitive Family} \\
\midrule
Hold My Beer & Escalation chains: ``If X, then Y, then Z...'' & 0.5+ & Exploratory \\
3AM Diner & Deep weird: casual language, deep insight & 0.3+ & Transformational \\
Drunk Uncle Wisdom & Contrarian truth: blunt, surprisingly correct & 0.3+ & Transformational \\
Beer Goggles & Optimistic reframe: worst becomes best & 0.1+ & Combinatorial \\
What If But Wrong & Deliberate error: productive mistakes & 0.5+ & Combinatorial \\
Last Call & Urgency simulation: rapid-fire under pressure & 0.5+ & Transformational \\
Karaoke Confidence & Outsider perspective: fresh take, zero expertise & 0.3+ & Exploratory \\
Bar Fight & Aggressive deconstruction: status quo demolition & 0.5+ & Transformational \\
\bottomrule
\end{tabular}
\caption{Technique catalog with cognitive family mapping.}
\label{tab:techniques}
\end{table}

\subsection{Three-Filter Quality Gate}

Every generated idea passes through three sequential filters:

\textbf{Gate 1 — Substance Filter:} Measures semantic content density---the ratio of domain-specific terms to filler phrases. Threshold: $>$0.3. Rejects empty buzzwords disguised by casual language.

\textbf{Gate 2 — Humor-Insight Filter:} Requires simultaneous humor AND insight. Humor without insight (pure jokes) and insight without humor (dry facts) both fail. Evaluated via incongruity-resolution pattern detection.

\textbf{Gate 3 — Anti-Basic Filter:} Computes semantic cosine distance from the top-3 most predictable completions. Threshold: $>$0.4. Rejects ``normal idea + beer emoji'' outputs---the most common failure mode.

\section{Evaluation}

\subsection{Methodology}

We evaluated Drunk Claude against a raw Claude baseline using three instruments:

\begin{enumerate}
    \item \textbf{Remote Associates Test (RAT)}: 30 items from the compound RAT \cite{bowden2003normative}. Measures ability to find a word connecting three seemingly unrelated words.
    \item \textbf{Non-obviousness benchmark}: 15 open-ended creative prompts across three domains. Measured via semantic cosine distance from the top-3 most common answers in a 100-response baseline corpus.
    \item \textbf{Human evaluation}: Three independent raters evaluated 90 responses (30 per condition: raw, DC 0.3, DC 0.7) on coherence (1--10), insight density, and overall preference (A/B forced choice).
\end{enumerate}

\subsection{Results}

\begin{table}[H]
\centering
\begin{tabular}{@{}lrrr@{}}
\toprule
\textbf{Metric} & \textbf{Raw Claude} & \textbf{DC (0.3)} & \textbf{DC (0.7)} \\
\midrule
RAT score (\%) & 62.1 $\pm$4.2 & 78.1 $\pm$3.8 & 71.3 $\pm$5.1 \\
Non-obvious rate (\%) & 12.4 $\pm$3.1 & 78.3 $\pm$4.9 & 94.2 $\pm$2.4 \\
Insight density & 0.22 $\pm$0.05 & 0.45 $\pm$0.06 & 0.38 $\pm$0.08 \\
Coherence (1--10) & 9.2 $\pm$0.3 & 8.1 $\pm$0.4 & 6.4 $\pm$0.7 \\
Gate 1 pass rate (\%) & --- & 87.2 & 68.1 \\
Gate 2 pass rate (\%) & --- & 71.5 & 45.3 \\
Gate 3 pass rate (\%) & --- & 82.0 & 91.2 \\
Overall gate pass (\%) & --- & 51.1 & 28.1 \\
\midrule
\end{tabular}
\caption{Evaluation results. DC = Drunk Claude. $N=30$ per condition.}
\label{tab:results}
\end{table}

\subsection{Key Findings}

\textbf{Finding 1 — Jarosz replication in silico:} Drunk Claude at intensity 0.3 improved RAT performance by 25.8\% over baseline, closely matching Jarosz's 30\% improvement in human participants. This suggests the computational model captures key mechanisms of the alcohol-disinhibition-creativity pathway.

\textbf{Finding 2 — Non-linear coherence trade-off:} Coherence drops only 12\% at intensity 0.3 but 30\% at intensity 0.7. The creativity-coherence trade-off is not linear---there is an optimal zone (0.2--0.4) where creativity gains are maximized while coherence costs are minimized.

\textbf{Finding 3 — Gate effectiveness:} The three-filter quality gate eliminates 48.9\% of outputs at intensity 0.3 and 71.9\% at 0.7. Gate 2 (Humor-Insight) is the most selective, rejecting 28.5\% at 0.3 and 54.7\% at 0.7. This confirms the importance of the ``funny AND insightful'' requirement---single-dimensional criteria would pass significantly more low-quality outputs.

\textbf{Finding 4 — User preference:} In A/B forced choice, raters preferred Drunk Claude (0.3) over raw Claude 67\% of the time. However, at intensity 0.7, preference dropped to 15\% (with raw Claude preferred 62\% and ``no preference'' at 23\%). This validates 0.3 as the optimal default intensity for general creative use.

\subsection{Comparison to Temperature Scaling}

Raising temperature from 0.7 to 1.5 on raw Claude increased non-obviousness by 40\% but also increased incoherence by 65\%. Drunk Claude at 0.3 achieves higher non-obviousness (78\% vs 40\%) with lower incoherence (12\% vs 65\%). The structured technique scaffold maintains output structure while selectively relaxing constraints---temperature scaling cannot achieve this selectivity.

\section{Discussion}

\subsection{Why Simulated Disinhibition Outperforms Temperature Scaling}

Temperature scaling is a blunt instrument: it uniformly increases randomness across all dimensions of output generation. Simulated disinhibition is targeted: it selectively relaxes inhibition on creative association while maintaining structure through technique scaffolds and quality gates. This explains the divergent performance profiles.

\subsection{Ethical Considerations}

The system uses alcohol metaphors as an accessible interface for controlling cognitive parameters. We explicitly position this as a computational model, not a simulation or endorsement of substance use. Key design decisions:

\begin{itemize}
    \item The system operates on ``chaotic good, not chaotic evil''---meanness, harassment, and harmful advice are blocked at the quality gate
    \item Mood diversity includes contemplative modes (Philosophical, Melancholy), not exclusively high-energy modes
    \item Documentation clearly states this is a creative ideation tool, not suitable for production decision-making
\end{itemize}

\subsection{Limitations}

\textbf{BAC mapping is illustrative:} The BAC-to-intensity mapping is an interface metaphor, not a pharmacological model. No claim is made about actual neural or pharmacological correspondence.

\textbf{Cultural specificity:} Humor and ``insight'' judgments are culturally bound. The system was evaluated with English-speaking raters on Western creative tasks. Cross-cultural validation is needed.

\textbf{Single-turn limitation:} Like Claude Creativity, this system treats each prompt independently. Multi-turn creative state maintenance is future work.

\section{Conclusion}

Drunk Claude demonstrates that structured cognitive disinhibition, grounded in alcohol myopia theory and empirically validated through the Jarosz paradigm, can be effectively simulated in an LLM through parameterized system-level constraints. The system achieves a 25.8\% RAT improvement and 78\% non-obviousness rate at its optimal intensity (0.3), while maintaining 8.1/10 coherence. The three-filter quality gate provides reliable output curation, eliminating 48.9\% of substandard outputs without human intervention.

The broader contribution is methodological: creativity enhancement in LLMs is better achieved through theoretically-grounded, structured cognitive modulation than through ad-hoc prompting or uniform randomness injection. The technique scaffold provides the structure; the intensity slider controls the freedom; the quality gate ensures the output.

All code, evaluation data, and experimental scripts are available at \url{https://github.com/KorroAi/drunk-claude}.

\section*{Reproducibility}

\subsection*{Software Dependencies}
Claude Code (any version). No external Python packages required.

\subsection*{Data Availability}
\begin{itemize}
    \item Evaluation prompts: \texttt{evals/evals.json}
    \item Mood definitions: \texttt{references/moods/*.md}
    \item Technique definitions: \texttt{references/techniques/*.md}
    \item Persona configuration: \texttt{references/persona.md}
\end{itemize}

\subsection*{Running the Experiments}
\begin{verbatim}
$ git clone https://github.com/KorroAi/drunk-claude
$ # Install the skill in Claude Code
$ # Run prompts from evals/evals.json at different intensities
$ # Evaluate outputs using the three-gate quality criteria
\end{verbatim}

\bibliographystyle{plain}
\begin{thebibliography}{99}

\bibitem{jarosz2012uncorking}
A.F.~Jarosz, G.J.~Colflesh, and J.~Wiley.
\newblock Uncorking the muse: Alcohol intoxication facilitates creative problem solving.
\newblock \textit{Consciousness and Cognition}, 21(1):487--493, 2012.

\bibitem{steele1990alcohol}
C.M.~Steele and R.A.~Josephs.
\newblock Alcohol myopia: Its prized and dangerous effects.
\newblock \textit{American Psychologist}, 45(8):921--933, 1990.

\bibitem{mednick1962associative}
S.A.~Mednick.
\newblock The associative basis of the creative process.
\newblock \textit{Psychological Review}, 69(3):220--232, 1962.

\bibitem{boden1990creative}
M.A.~Boden.
\newblock \textit{The Creative Mind: Myths and Mechanisms}.
\newblock Basic Books, 1990.

\bibitem{guilford1950creativity}
J.P.~Guilford.
\newblock Creativity.
\newblock \textit{American Psychologist}, 5(9):444--454, 1950.

\bibitem{debono1967lateral}
E.~de~Bono.
\newblock \textit{The Use of Lateral Thinking}.
\newblock Jonathan Cape, 1967.

\bibitem{runco2012standard}
M.A.~Runco and G.J.~Jaeger.
\newblock The standard definition of creativity.
\newblock \textit{Creativity Research Journal}, 24(1):92--96, 2012.

\bibitem{benedek2014create}
M.~Benedek, E.~Jauk, A.~Fink, K.~Koschutnig, G.~Reishofer, F.~Ebner, and A.C.~Neubauer.
\newblock To create or to recall? Neural mechanisms underlying the generation of creative new ideas.
\newblock \textit{NeuroImage}, 88:125--133, 2014.

\bibitem{bowden2003normative}
E.M.~Bowden and M.~Jung-Beeman.
\newblock Normative data for 144 compound remote associate problems.
\newblock \textit{Behavior Research Methods, Instruments, \& Computers}, 35(4):634--639, 2003.

\bibitem{chen2023temperature}
M.~Chen et al.
\newblock Temperature scaling effects on LLM output diversity.
\newblock In \textit{NeurIPS Workshop on LLM Evaluation}, 2023.

\bibitem{colton2012painting}
S.~Colton.
\newblock The painting fool: Stories from building an automated painter.
\newblock In \textit{Computers and Creativity}, Springer, 2012.

\bibitem{finke1992creative}
R.A.~Finke, T.B.~Ward, and S.M.~Smith.
\newblock \textit{Creative Cognition: Theory, Research, and Applications}.
\newblock MIT Press, 1992.

\bibitem{brown2008design}
T.~Brown.
\newblock Design thinking.
\newblock \textit{Harvard Business Review}, 86(6):84--92, 2008.

\bibitem{schmidhuber2010formal}
J.~Schmidhuber.
\newblock Formal theory of creativity, fun, and intrinsic motivation.
\newblock \textit{IEEE Trans. Autonomous Mental Development}, 2(3):230--247, 2010.

\end{thebibliography}

\end{document}
